% ======================================================================
%  Section 4.8 --- Conclusion.  WRITTEN.
%
%  FOUR DECISIONS I MADE IN WRITING IT, each of which you may want to
%  reverse, and one of which you should think about before you do.
%
%  1. IT CLAIMS ONE RESULT, NOT TWO.
%     The chapter has a decisive result -- robust weighting buys two to
%     three orders of magnitude under occlusion -- and an undecided one --
%     whether the structure-informed measure improves on generic Tukey
%     weighting. 4.7.2 reports the second as a tie leaning to Variant 2 in
%     translation. A conclusion that presents both as achievements will be
%     read as overclaiming and will cost you the first as well as the
%     second. So the conclusion claims the first, states the second as
%     undecided, and rests the case for the proposed weighting on the
%     formulation and the mechanism, where the evidence actually is.
%     This is not a weak conclusion. "Robust weighting is necessary in
%     direct visual servoing and here is a four-hundred-fold demonstration"
%     is a strong chapter. It just is not the chapter you set out to write,
%     and it is better to say so than to be asked.
%
%  2. IT REPORTS THE GAIN-SCHEDULE FINDING AS A CONTRIBUTION.
%     It was arrived at accidentally, by way of explaining why the two
%     tables disagree by three orders of magnitude. It is nonetheless the
%     most actionable thing in the chapter and it is quantified to 0.62 per
%     cent. Burying it in 4.7.2 would waste it.
%
%  3. IT REPORTS THE DC-FREE CORRECTION AS A METHODOLOGICAL FINDING.
%     That summed Gauss-Hermite kernels have a DC gain equal to their L1
%     norm, so that a structure measure built from them reports local
%     brightness and not local structure, is a transferable observation
%     that anyone building such a measure would want. One sentence.
%
%  4. IT DOES NOT SUMMARISE THE CHAPTER SECTION BY SECTION.
%     A conclusion that recapitulates 4.2, then 4.3, then 4.4 is the part
%     of a thesis nobody reads. This one states what is now known that was
%     not known at the start of the chapter, which is a different and
%     shorter list.
%
%  NOTE ON NUMBERS: every figure quoted here appears in 4.7, at iteration
%  600, in the corrected units. No new quantity is introduced. If you
%  change a number in 4.7, it changes here -- there are seven of them and
%  they are listed at the foot.
% ======================================================================

\section{Conclusion}
\label{sec:ch4-conclusion}

The photometric scheme of Chapter~\ref{ch:herpvs} uses every pixel of the
image as a measurement, and takes its accuracy from that redundancy. The same
redundancy is its weakness: the control law assumes that every measurement is
explained by the pose, and an occlusion, a specular highlight or a moving
object violates that assumption over a compact region of the image while
leaving the rest of it intact. This chapter addressed that vulnerability, and
did so by treating the difficulty as one of estimation rather than of
optimisation. The cost is not hard to minimise in the presence of an occluder;
minimising it is the wrong thing to do, because the occluded measurements do
not raise the cost at the correct pose but displace the pose at which the cost
is least.

Two things were proposed. The first is the weighted control
law~\eqref{eq:robust-lm-law}, in which a diagonal weighting acts on the
feature error and the interaction matrix alike, with the weights recomputed at
every iteration from a redescending estimator and a robust scale. Setting the
weights to unity recovers the scheme of Chapter~\ref{ch:herpvs} up to the
floor~\eqref{eq:hessian-floor} on the damped Hessian, so the unweighted scheme
is the degenerate case of the formulation rather than a separate method placed
beside it. The second is the structure-informed weight
of~\eqref{eq:hermite-weights}, which normalises each residual by a
scale-aggregated measure of the local Hermite structure before submitting it
to the estimator. Its justification is not heuristic: by the factorisation
established in Section~\ref{sec:hermite-principle} the norm of a row of the
interaction matrix is the product of the local gradient magnitude and a term
depending only on the image coordinates and the depth, and the structure
measure is a scale-aggregated gradient magnitude, so a measurement carrying
little structure also constrains the pose weakly, whatever the cause of its
flatness.

The principal result is quantitative and is not in doubt. Under an occlusion
covering some six per cent of the measurements, the unweighted scheme settles
$1293\,\mu\mathrm{m}$ and $1.70^{\circ}$ from the desired pose, while both
weighted variants settle within $3.2\,\mu\mathrm{m}$ and $0.017^{\circ}$: a
factor of between four and five hundred in translation and of about a hundred
in rotation. Robust weighting is therefore not a refinement of a direct scheme
but a condition of its usefulness outside the laboratory.

The manner of the unweighted scheme's failure is the more instructive half of
that result, and it is the clearest demonstration in this thesis of why the
estimation framing is the right one. Variant~1 did not diverge, oscillate or
stall prematurely, and it did not come to rest at a higher cost than the
weighted variants. It attained the \emph{smallest} final feature error of the
three, lower by $1.4\%$, at a pose error larger by a factor of four hundred,
and its pose was unchanged to six significant figures over the last two
hundred iterations. It converged accurately upon the wrong answer. The
consequence for practice is immediate: no criterion computed from the same
corrupted measurements that drive the control can detect this failure, and a
practitioner monitoring the residual would have judged the unweighted run the
most successful of the three.

Against that, the chapter also measured what weighting costs when there is
nothing to reject, which the literature more often assumes than reports. On
the clean target the weighted variants are the less accurate, by a few tenths
of a degree in rotation and a fraction of a millimetre in translation, and the
degradation falls almost entirely on the depth and the rotation coupled to it
--- the two degrees of freedom that a fronto-parallel photometric measurement
constrains least securely. The weight maps of Section~\ref{sec:weight-maps}
show why the penalty does not disappear at convergence: the scale is estimated
from the residuals and contracts with them, so the mean weight on a clean
scene at convergence is $0.789$ and only one measurement in seven carries
full weight. An M-estimator with an adaptive scale does not switch itself off
when the data cease to contain outliers, and the cost of robustness on a clean
scene is accordingly permanent rather than transient. It is a cost worth
paying at the exchange rate the occlusion experiment establishes, but it
should be paid knowingly.

The weight maps also settled how the mechanism operates, and the answer was
not the expected one. The rejection of the occluder is not available at the
outset and does not enable the convergence; it is produced by the convergence.
At the first iteration the misalignment residual dominates the image, the
robust scale is large in proportion, and at most a tenth of the occluded
region is rejected --- against a third of the clean scene's valid measurements
carried below a quarter weight, which is the over-rejection that motivated the
structure normalisation in the first place and which the normalisation reduces
but does not remove. Only as the pose error falls, and the scale with it, does
the occluder become the sole large residual in the image and is rejected
decisively, as a solid region covering the measurements that are corrupted.
Robust weighting and pose convergence reinforce one another.

Three further findings should be recorded, none of them anticipated when the
chapter was begun. The first is that a structure measure built from summed
Gauss--Hermite kernels cannot serve, because their gain at zero frequency
equals their $L_{1}$ norm: such a measure reports local \emph{brightness} and
not local structure, and the measure must be built from the odd, zero-mean
kernels for the argument of Section~\ref{sec:hermite-principle} to hold of it
at all. The second is that the fixed floor~\eqref{eq:hessian-floor} renders
the rank condition of Section~\ref{sec:weighted-law} non-binding by
construction, so that no monitoring of the conditioning is required; what the
floor cannot restore is the information the rejected measurements carried, and
the consequence of rejecting too many is a loss of accuracy rather than a
divergence. The third is the most useful of the three and was arrived at
indirectly. Because the adaptive gain is proportional to the intensity error,
and the commanded velocity to the residual in addition, the velocity decays
roughly as the cube of the error and the nominal runs stall while a
sub-millimetre misalignment remains. Under occlusion the intensity error
cannot decay, a fixed residual persisting wherever the reference has no
counterpart, so the gain is held some fifty-four times higher, the cubic
collapse is replaced by a velocity linear in the error, and the final accuracy
improves by three orders of magnitude. The relation is not an interpretation
placed on the data: since the damping depends on the same intensity error with
a fifth-power relation to the gain, the damping ratio can be predicted from
the gain ratio, and across three runs the prediction agrees with the
measurement to within $0.62\%$, the exponent recovered from the data being
$9.997$ against the $10$ of~\eqref{eq:mu}. The stall of the nominal runs is
therefore a property of the gain schedule inherited from the implementation
and not of the weighting or of the features, and a gain that does not vanish
with the error is the single most valuable improvement the chapter identifies.

What the chapter does not establish should be stated as plainly. It does not
rank the three variants. Each figure rests on one run from one initial pose,
and the two weighted variants are not separated by the occlusion experiment:
Tukey weighting attained the smaller translational error, by $28\%$, and the
structure-informed weighting the smaller rotational error, by $8\%$, at
magnitudes --- under a micrometre apart --- at which the difference is without
practical consequence and cannot be shown to be systematic. The proposed
weighting is therefore not demonstrated here to position more accurately than
a generic robust estimator; what is demonstrated is that it reaches comparable
accuracy by a different route, suppressing measurements for what the image
lacks rather than for how far they depart from a pose estimate that is itself
in error. The one systematic difference observed between the two is not in the
proposal's favour: the structure-informed variant comes to rest with a
terminal jitter thirty-four times larger, its weights continuing to vary after
its pose has ceased to improve. Nor was a weight map recorded for the
residual-based variant, so even the claim about the mechanism is not
comparative. Establishing any of this would require a study over several
initial poses, together with a weight map for the residual-based variant and a
direct measurement of the overhead of the weighting, none of which is
undertaken here.

The contribution of the chapter is accordingly narrower than a ranking and
firmer than one. A robust, structure-informed weighting has been formulated
for the Hermite-based photometric scheme, given a justification that rests on
the structure of the interaction matrix rather than on intuition, situated
within the stability analysis of the unweighted case, and evaluated in an
ablation in which the weighting is the only quantity varied. Within that
ablation, robust weighting is shown to be decisive against occlusion, the
failure it prevents is shown to be undetectable from the cost, the price of
robustness on uncorrupted data is measured rather than assumed, and the
mechanism by which the weighting acts is made directly observable. Whether the
structural information improves upon a generic estimator remains open, and the
experiments needed to close it are named.

% ======================================================================
%  THE SEVEN NUMBERS QUOTED HERE, all from 4.7 at iteration 600. If any
%  changes there, change it here.
%    1293 um, 1.70 deg          Variant 1 occluded      (4.7.2)
%    3.2 um, 0.017 deg          bound on Variants 2, 3  (4.7.2)
%    1.4 per cent               V1's feature error deficit (4.7.2)
%    0.789                      mean weight, clean, 600 (4.7.3)
%    28 per cent, 8 per cent    V2/V3 margins           (4.7.2)
%    54, 0.62 per cent, 9.997   gain and damping        (4.7.2)
%    thirty-four                V3 terminal jitter      (4.7.2)
%  Two proportions are quoted in words rather than figures -- "at most a
%  tenth of the occluded region" and "a third of the clean scene's valid
%  measurements" (13.6 per cent) -- because a conclusion reads badly when
%  it is dense with decimals. Both are exact as stated.
%
%  LABELS USED
%    ch:herpvs                Chapter 3
%    eq:robust-lm-law         the weighted control law (4.4)
%    eq:hessian-floor         the floor on the damped Hessian (4.2.2)
%    eq:hermite-weights       the structure-informed weights (4.13)
%    eq:mu                    the damping
%    sec:hermite-principle    4.4.1
%    sec:weighted-law         4.2.2
%    sec:weight-maps          4.7.3
%    sec:future               future work
%
%  ONE THING TO DECIDE. The seventh paragraph is long -- three findings in
%  one paragraph, the third of them substantial. If you would rather, the
%  gain-schedule finding can stand alone as its own paragraph or even as a
%  short numbered subsection, since it is the chapter's most actionable
%  result and a reader skimming the conclusion should not be able to miss
%  it. I have left it inside the paragraph so that the conclusion does not
%  appear to end on a result about the implementation rather than about
%  the method; move it if you disagree.
%
%  AND ONE THING I DELIBERATELY LEFT OUT. The computational cost. The
%  timings in the logs measure machine load, not the algorithm -- Variant 3
%  runs fastest of the three in one set and slowest in another -- so no
%  cost claim can be made from them, and a conclusion is the wrong place
%  to introduce a caveat of that kind. If you instrument the weight
%  computation as suggested in 4.7.2, a sentence here would be worth
%  having: the argument of 4.4.1 predicts about twenty per cent, and a
%  measured twenty per cent would close the chapter well.
% ======================================================================
